% Generated draft: requires mathematical and source review.
\documentclass[11pt]{article}
\usepackage[margin=25mm]{geometry}
\usepackage{fontspec}
\setmainfont{Latin Modern Roman}
\usepackage{amsmath,amssymb,mathtools,unicode-math}
\setmathfont{Latin Modern Math}
\usepackage{enumitem,needspace,xurl,hyperref}
\setlength{\parindent}{0pt}
\setlength{\parskip}{5pt}
\setlength{\emergencystretch}{3em}
\newcommand{\N}{\mathbb{N}}
\newcommand{\Z}{\mathbb{Z}}
\newcommand{\Q}{\mathbb{Q}}
\newcommand{\R}{\mathbb{R}}
\newcommand{\C}{\mathbb{C}}
\newcommand{\F}{\mathbb{F}}
\newcommand{\A}{\mathbb{A}}
\newcommand{\PP}{\mathbb{P}}
\newcommand{\E}{\mathbb{E}}
\newcommand{\eps}{\varepsilon}
\newcommand{\dd}{\,\mathrm{d}}
\begin{document}
\section*{UnsolvedMath ID: 30002068}
\textbf{Source problem code:} catalogue entry
\subsection{Definitions and mathematical statement}
Let \(d\geq 1\). A lattice polytope in \(\R^d\) is the convex hull of finitely many points of \(\Z^d\). A full-dimensional lattice polytope is smooth if exactly \(d\) edges meet at every vertex and their primitive integral direction vectors pointing away from that vertex form a basis of \(\Z^d\). A full-dimensional lattice polytope containing the origin in its interior is reflexive if its polar \(P^\vee=\{y\in\R^d:\langle y,x\rangle\leq 1\text{ for every }x\in P\}\) is a lattice polytope.

A lattice triangulation of a lattice polytope \(P\) is a finite simplicial complex with lattice vertices whose union is \(P\), with any two simplices intersecting in a common face or not at all. A full-dimensional lattice simplex \(\operatorname{conv}(v_0,\ldots,v_d)\) is unimodular if \(\lvert\det(v_1-v_0,\ldots,v_d-v_0)\rvert=1\). A triangulation is unimodular if all its maximal simplices are unimodular. It is regular if there is a height function \(h:P\cap\Z^d\to\R\) such that the projections of the lower faces of \(\operatorname{conv}\{(a,h(a)):a\in P\cap\Z^d\}\) give that triangulation. Polytopes of smaller affine dimension can be treated in their affine lattice.

\textbf{General question.} Does every smooth lattice polytope, in every dimension, admit a unimodular lattice triangulation?

\textbf{Specified example.} Define the set given by the supplied inequalities by
\[
P=\left\{x\in\R^7:\begin{aligned}
x_2-2x_6-x_7&\leq 1,\\
2x_6-x_1-x_2&\leq 1,\\
2x_1-x_2-x_3-x_4-x_5+x_6&\leq 1,\\
x_2-x_6&\leq 1,\\
x_i&\leq 1\qquad(i=1,\ldots,7)
\end{aligned}\right\}.
\]
The source describes \(P\) as a smooth reflexive \(7\)-dimensional lattice polytope and these inequalities as its facets; those descriptions are retained as supplied qualifications, not independently established here. Does \(P\) admit a regular unimodular lattice triangulation?

For a field \(k\), put \(A=P\cap\Z^7\). Define the toric ideal of the lattice-point configuration of \(P\) by
\[
I_{P,k}=\ker\!\left(k[z_a:a\in A]\longrightarrow k[s,t_1^{\pm1},\ldots,t_7^{\pm1}],\quad z_a\longmapsto s\prod_{i=1}^7t_i^{a_i}\right).
\]
Give each \(z_a\) degree one. Is \(I_{P,k}\) generated by homogeneous quadratic binomials? The source does not specify the ground field or its quantification, nor explicitly define its toric-ideal convention. The displayed definition uses all lattice points of \(P\) and records the usual homogeneous convention; the intended convention and field scope require confirmation.

\textbf{Unresolved source qualification.} The supplied dated literature assessment instead discusses an incomplete heat-kernel problem and does not address these polytope questions. It therefore cannot clarify their formulation or substantiate their present status. The supplied labels \texttt{open} and \texttt{corrected\_verified} are registry annotations, not independent mathematical verification.
\subsection{Short English statement}
The record asks whether every smooth lattice polytope has a unimodular triangulation. It also asks whether the explicitly specified \(7\)-dimensional polytope has a regular unimodular triangulation and whether its toric ideal is generated by quadrics. The toric-ideal convention and ground-field scope need confirmation, and the supplied literature assessment concerns an unrelated heat-kernel problem.
\subsection{Sources}
\begin{itemize}
\item Supplied local problem registry: UnsolvedMath permanent ID \(30002068\), source code \texttt{catalogue entry}, title \emph{Unimodular Triangulations of Smooth Lattice Polytopes}, in \emph{Oberwolfach Reports Open Problems}. The proposed year is recorded as 2012; no proposer is named.
\item Supplied source citation: \emph{Triangulations (2013)}, Oberwolfach Reports, pp.~1405--1486, DOI \texttt{10.4171/owr/2012/24}, \url{https://doi.org/10.4171/owr/2012/24}. This is the record's report attribution; its correspondence to the problem has not been independently verified.
\item Supplied literature source, labeled ``Lesch analytic-torsion paper'': \url{https://doi.org/10.2140/apde.2013.6.221}.
\item Supplied literature source, labeled ``original OWR report'': \url{https://doi.org/10.4171/owr/2012/25}. This DOI differs from the main source citation.
\item Literature assessment dated exactly \texttt{2026-08-21}: the supplied curation labels the problem open but discusses an incomplete heat-kernel formulation. This is supplied dated curation, not an independently verified literature assessment. The background expressly requests independent verification of time-sensitive status and cited claims.
\end{itemize}
\end{document}
